\documentclass[11pt,letterpaper]{report}

\usepackage[utf8]{inputenc}
\usepackage[T1]{fontenc}
\usepackage[spanish,provide=*]{babel}
\usepackage{amsmath,amssymb}
\usepackage{booktabs}
\usepackage{array}
\usepackage{graphicx}
\usepackage{float}
\usepackage{caption}
\usepackage[margin=2.5cm]{geometry}
\usepackage{enumitem}

\renewcommand{\arraystretch}{1.15}

\begin{document}

\begin{center}
{\LARGE\bfseries\linespread{1.3}\selectfont 7. LENGUAJES DE SIMULACIÓN\par}
\end{center}
\vspace{0.5cm}

\setcounter{chapter}{7}
\setcounter{page}{123}

%==========================================================
Las primeras etapas de un estudio de simulación se refieren a la
definición del sistema a ser modelado y a la descripción del sistema de
términos de relaciones lógicas de sus variables y diagramas de flujo.
Sin embargo, llega el momento de describir el modelo en un lenguaje que
sea aceptado por la computadora que se va a usar. En esta etapa se
tienen dos cursos de acción a seguir: 1) Desarrollar el software
requerido para estudios de simulación, o 2) Comprar algún paquete de
simulación como GPSS, GASP, simula, etc. Desafortunadamente, para la
segunda alternativa existen una gran cantidad de paquetes (lenguajes de
programación de propósito especial), que resulta extremadamente difícil
decidir cuál paquete se ajusta mejor a una aplicación determinada. Esta
situación origina que en la mayoría de las veces, la selección de un
paquete depende de si el analista lo conoce, lo entiende y lo sabe
aplicar. No obstante, estas dos alternativas se deben de evaluar en
términos económicos antes de tomar una decisión.

%==========================================================
\section*{7.1. Ventajas de los lenguajes de simulación}

El proceso evolutivo de los lenguajes de simulación ha sido largo y
extenso. Empezó a finales de la década de los 50's. En un principio los
lenguajes que se desarrollaron eran de propósito general. Sin embargo,
poco a poco los estudiosos de este tema se dieron cuenta de la gran
similitud que existía entre las diferentes situaciones o sistemas que se
simulaban. Lo anterior condujo obviamente al desarrollo de lenguajes de
propósito especial, los cuales en la actualidad tienen una gran demanda
y su proceso de comercialización ha sido amplio y extenso. Entre las
ventajas principales de estos lenguajes de simulación, se pueden
mencionar las siguientes:

\begin{enumerate}[label=\alph*)]
\item Reducción en la tarea de programación. Con los lenguajes de
simulación, el tiempo dedicado a la programación del modelo se reduce
considerablemente. Existen algunos paquetes como GPSS, en los que con un
número muy reducido de estatutos, se pueden simular sistemas que en otro
lenguaje como Fortran, requerirían una gran cantidad de estatutos y
subrutinas.
\item Mejor definición del sistema. A través de los lenguajes de
simulación, se facilita la tarea de definir las diferentes entidades que
interactúan dentro del sistema. También, con estos lenguajes se
determina con mayor facilidad las interrelaciones que existen entre las
entidades que forman el sistema.
\item Mayor flexibilidad para cambios. Con los lenguajes generales como
Fortran, el proceso de cambios puede ser largo y tedioso. Sin embargo,
con el uso de lenguajes de simulación, los cambios son una tarea simple
y rutinaria.
\item Mejor diferenciación de las entidades que forman el sistema. El
uso de lenguajes de simulación facilita determinar o definir las
características y atributos de una entidad. Con las entidades bien
definidas y diferenciadas, se aumenta y mejora el entendimiento del
sistema a simular.
\item Se relacionan mejor las entidades. Con las entidades bien
definidas, los lenguajes de simulación permiten relacionar mejor a cada
una de estas entidades, es decir, se determina más fácilmente las
relaciones que las entidades guardan entre sí y el análisis de cada una
de ellas.
\end{enumerate}

%==========================================================
\section*{7.2. Características de los lenguajes de simulación}

Los lenguajes de simulación que actualmente existen en el mercado,
tienen una serie de características propias que los diferencian de los
demás. Entre estas características se pueden mencionar las siguientes:

\begin{enumerate}[label=\alph*)]
\item El procedimiento utilizado para generar números aleatorios
uniformes.
\item Los procedimientos o métodos utilizados para generar las
variables aleatorias no-uniformes más conocidas y más usadas.
\item La forma de adelantar el ``reloj de la simulación'', la cual
puede ser: 1) Incrementos a tiempo fijo, o 2) Incrementos al próximo
evento.
\item El análisis estadístico de los resultados de la simulación.
\item El formato en que los resultados de la simulación son
presentados.
\item La forma en que las inconsistencias y errores de lógica es
reportada.
\item El lenguaje en el cual el paquete está escrito, el cual puede
ser: Fortran, Algol, PL/I, Asembler, etc.
\item Los diferentes tipos de computadoras cuyo compilador es
compatible con el del paquete en cuestión.
\end{enumerate}

A continuación, se presentan las características principales de los
lenguajes de simulación más usados.

\medskip
\noindent\textbf{GPSS (General Purpose Simulation System)}

\begin{table}[H]
\centering
\begin{tabular}{ll}
Persona que lo desarrolló: & Geoffrey Gordon \\
Versiones más conocidas: & GPSS I, GPSS II, GPSS III, GPSS/360, GPSS V. \\
Lenguaje del paquete: & Asembler \\
Reloj de la simulación: & Incrementos al próximo evento \\
Computadoras compatibles: & Generalmente se adapta a cualquier tipo de computadora. \\
\end{tabular}
\end{table}

\noindent\textbf{SIMSCRIPT (No tiene ningún significado)}

\begin{table}[H]
\centering
\begin{tabular}{ll}
Personas que lo desarrollaron: & H. M. Markowitz, H. W. Karr y B. Hausner. \\
Versiones más conocidas: & Simscript I, Simscript I.5, Simscript II, Simscript II.5, C-Simscript. \\
Lenguaje del paquete: & Fortran las primeras versiones, Asembler las últimas. \\
Reloj de la simulación: & Incrementos al próximo evento para el caso discreto, e incrementos a tiempo fijo para el caso continuo (C-Simscript). \\
Computadoras compatibles: & CDC 6000/7000, univac 1100, IBM 360/370, Honeywell 600/6000. \\
\end{tabular}
\end{table}

\noindent\textbf{GASP (General Activity Simulation Program)}

\begin{table}[H]
\centering
\begin{tabular}{ll}
Personas que lo desarrollaron: & P. J. Kiviat y A. Colher \\
Versiones más conocidas: & GASP II, GASP IV, GASP-PLUS \\
Lenguaje del paquete: & Fortran, PL/I. \\
Reloj de la simulación: & Incrementos al próximo evento para el caso discreto, e incrementos a tiempo fijo para el caso continuo (GASP IV y PLUS). \\
Computadoras compatibles: & Cualquier computadora con compilador de Fortran o PL/I. \\
\end{tabular}
\end{table}

\noindent\textbf{SLAM (Simulation Language for Alternative Modeling)}

\begin{table}[H]
\centering
\begin{tabular}{ll}
Personas que lo desarrollaron: & A. Alan B. Pritsker y Asociados \\
Versiones más conocidas: & SLAM fue el resultado de la fusión de varios lenguajes como GASP IV y QGERT. \\
Lenguaje del paquete: & Fortran IV. \\
Reloj de la simulación: & Incrementos al próximo evento para el caso discreto, e incrementos a tiempo fijo para el caso continuo. \\
Computadoras compatibles: & Cualquier computadora con compilador de Fortran \\
\end{tabular}
\end{table}

Cualesquiera de estos lenguajes tienen sus propias ventajas y
desventajas y no se puede decir que un lenguaje es mejor que otro.
Generalmente, entre más fácil de aprender y de usar sea un lenguaje,
menor será su flexibilidad y su eficiencia. Por consiguiente, decidir
qué lenguaje utilizar en una aplicación específica, no es una tarea
fácil de realizar.

%==========================================================
\section*{7.3. Factores a considerar en la selección de un lenguaje}

La selección de un lenguaje de simulación generalmente está supeditada
al tipo de computadora que se tiene disponible, es decir, en la mayoría
de las veces ya se cuenta con cierta configuración de hardware. Por
consiguiente, conociendo la computadora que se va a usar, los factores a
considerar en la selección del lenguaje serían:

\begin{enumerate}[label=\alph*)]
\item Los manuales disponibles. Es muy importante considerar la
facilidad de entender e interpretar los manuales disponibles.
\item Compilador del lenguaje. Es necesario investigar la
compatibilidad del compilador del lenguaje con la computadora
disponible.
\item La documentación y diagnóstico de errores. Es conveniente
analizar la forma en que el lenguaje reporta las inconsistencias y los
errores de lógica.
\item La eficiencia. Uno de los factores principales a considerar en
la selección de un lenguaje es su eficiencia de operación. Dentro de la
eficiencia se considera el tiempo de organizar, programar, compilar y
ejecutar.
\item Los costos involucrados. Entre los costos que origina la
adquisición de un paquete se pueden mencionar: El costo de instalación,
el costo de mantenimiento y actualización del software y el costo de
operación.
\item Conocimiento del lenguaje. Otro factor importante a considerar
en la selección del lenguaje, es el conocimiento y dominio que de éste
tengan las personas o analistas encargados de realizar los estudios de
simulación.
\item Justificación económica. Finalmente, y el más importante de
todos, es la justificación económica del lenguaje de simulación. A este
respecto, es conveniente señalar que la adquisición de un paquete se
debe considerar como un proyecto de inversión, donde los beneficios que
se derivan de tal adquisición, deben de compensar la inversión y los
gastos que origina.
\end{enumerate}

%==========================================================
\section*{7.4. Clasificación de los lenguajes de simulación}

Los modelos de simulación se clasifican en dos categorías: 1) Modelos de
simulación continua y 2) Modelos de simulación discreta. Los modelos de
simulación continua, generalmente son apropiados cuando el analista
considera al sistema bajo estudio, como un flujo continuo de
información. En simulación continua, generalmente el reloj de
simulación se incrementa a intervalos fijos de tiempo. Por otra parte,
en modelos de simulación discreta, al analista le interesa lo que le
sucede a entidades individuales del sistema. Por consiguiente, el reloj
de la simulación en simulación discreta, se incrementa cada vez que
ocurre un evento.

Los modelos de simulación discreta pueden desarrollarse a través de tres
enfoques: 1) Enfoque de eventos, 2) Enfoque de actividades y 3) Enfoque
de procesos. Un evento se define como un cambio en el estado de una
entidad del sistema. Una actividad es una colección de operaciones que
transforman el estado de una entidad. Un proceso es una secuencia de
eventos que ocurren en un tiempo determinado.

\begin{figure}[H]
\centering
\caption*{\textbf{Figura 7.1} Eventos, actividades y procesos}
\fbox{\parbox{0.85\textwidth}{\centering\vspace{2cm}
[Figura: línea de tiempo "Proceso" con Evento 1 (arribo), Evento 2 (empieza el servicio), Evento 3 (empieza el servicio), Evento 4 (termina el servicio), Evento 5 (termina el servicio); Actividad 1 entre Eventos 2-3, Actividad 2 entre Eventos 3-5]
\vspace{2cm}}}
\end{figure}

Además de los tres enfoques anteriores, algunos autores hacen una
diferenciación entre lenguajes con orientación simbólica (GPSS) y
lenguajes con orientación de estatutos de programación. Los primeros son
más fáciles de aprender, pero los segundos son más flexibles. La figura
7.2 muestra una clasificación de los lenguajes de simulación. En esta
figura se puede apreciar que los lenguajes simbólicos a pesar de tener
un enfoque de procesos en su diseño, se consideran como una cuarta
categoría para diferenciarlos de los lenguajes de orientación de
estatutos de programación.

Finalmente, vale la pena mencionar que recientemente se han estado
desarrollando lenguajes que permiten el análisis y evaluación de modelos
discretos y modelos continuos. GASP IV y C-SIMSCRIPT fueron los primeros
lenguajes de este tipo que se desarrollaron, aunque recientemente está
disponible un nuevo lenguaje de este tipo llamado SLAM. Estos lenguajes
permiten al analista modelar sistemas discretos, sistemas continuos y
sistemas donde algunas entidades se comportan en forma discreta y otras
en forma continua.

\begin{figure}[H]
\centering
\caption*{\textbf{Figura 7.2} Clasificación de los lenguajes de simulación}
\fbox{\parbox{0.85\textwidth}{\centering\vspace{2cm}
[Figura: árbol -- Lenguajes de simulación se divide en: Lenguajes para simulación continua / Lenguajes para simulación continua y discreta (GASP IV, C-SIMSCRIPT, SLAM) / Lenguajes para simulación discreta, que a su vez se divide en Enfoque de eventos (GASP II, SIMSCRIPT), Enfoque de actividades (FORSIM-IV), Enfoque de procesos (SIMULA), Enfoque de símbolos (GPSS)]
\vspace{2cm}}}
\end{figure}

%==========================================================
\section*{7.5. Introducción al GPSS}

Como se ha mencionado en párrafos anteriores, existen actualmente en el
mercado una gran cantidad de paquetes de simulación. Sin embargo, el
paquete GPSS es sin duda uno de los más conocidos y usados en empresas e
instituciones educativas.\footnote{Actualmente en el ITESM se está
utilizando la versión GPSS-V.} Por esta razón, en esta sección se
tratará de describir a través de una serie de ejemplos, las
potencialidades y los diferentes campos de aplicación en los cuales se
puede aplicar este paquete (versión GPSS-V). Obviamente esta
presentación no incluye toda la gama de aplicaciones que se le pueden
dar a este paquete. Para tener una concepción más amplia del GPSS, se
puede recurrir al libro \emph{Simulation Using GPSS} cuyo autor es
Thomas J. Schriber.

Para representar la llegada y salida de una facilidad de servicio
(1 servidor), los bloques que se utilizan son SEIZE y RELEASE, donde el
campo A representa la facilidad de servicio que se está utilizando. El
primer bloque SEIZE indica que se inicia la entrada a la facilidad de
servicio A y el bloque RELEASE representa la salida de dicha facilidad.

El tiempo de servicio se representa mediante el bloque ADVANCE A,B
donde: Campo A = Tiempo promedio de servicio y Campo B = Desviación con
respecto a la media. Por ejemplo si el tiempo entre llegadas es
uniforme entre 10 y 18 minutos, entonces, A = 14 y B = 4.

Finalmente, para llevar un conteo de la cantidad de transacciones
(clientes) que han pasado por el sistema, se utiliza el bloque
TERMINATE, y para representar la generación de llegadas al sistema se
utiliza el bloque GENERATE A,B,C,D,E donde: Campo C especifica el tiempo
de ocurrencia de la primera transacción, Campo D especifica el límite de
transacciones que pueden ocurrir en el bloque GENERATE, y Campo E
especifica la prioridad de las transacciones al momento de generarse.

\subsection*{Ejemplo 7.2. Centro Electrónico de Cálculo}

Al salón de terminales del Centro Electrónico de Cálculo del ITESM,
arriban estudiantes de acuerdo a una distribución uniforme entre 2 y 8
minutos. En el salón existen 10 terminales. El tiempo promedio que un
estudiante usa la terminal, está uniformemente distribuido entre 40 y
80 minutos. Si al llegar al salón de terminales el estudiante las
encuentra todas ocupadas, entonces, existe un 50\% de probabilidad de
que el estudiante regrese después de 5 minutos (suponga que en esos 5
minutos el estudiante va a la cafetería a comprar golosinas y
refrescos). ¿Determine al momento de que 1000 estudiantes hayan pasado
por el salón de terminales, la cantidad de estudiantes que al llegar al
sistema encontraron todas las terminales ocupadas y no regresaron?

Para simular este problema se requiere utilizar nuevos bloques. Por
ejemplo, es necesario conocer qué bloques utilizar para representar a
una facilidad de servicio con muchos servidores (cada terminal es un
servidor con un tiempo de servicio uniforme entre 40 y 80 minutos).
También, se requiere conocer el bloque que representa las
transferencias condicionales e incondicionales que ocurren en el
sistema. Para el primer caso, los bloques que se utilizan son ENTER y
LEAVE, donde: $A$ = Representa el nombre o el número que se le da a la
facilidad de servicio (STORAGE), y $B$ = Representa la cantidad de
servidores que se ocupan o desocupan a la vez, al llegar o abandonar el
sistema una determinada transacción. El primer bloque ENTER indica que
se inicia la entrada a la facilidad de servicio STORAGE A y el bloque
LEAVE representa la salida de dicha facilidad.

Por otra parte, para representar las transferencias condicionales e
incondicionales que ocurren en un sistema, se utiliza el bloque
TRANSFER, el cual puede ser codificado de las siguientes formas:

\begin{verbatim}
TRANSFER  .3, ETI 1, ETI 2
TRANSFER  BOTH,, ETI 1
TRANSFER  , ETI 2
\end{verbatim}

El primer caso indica que la probabilidad de ir al bloque con etiqueta
ETI 1 es de 0.70. El complemento es la probabilidad de ir al bloque con
etiqueta ETI 2. El segundo caso indica que la transacción continúa en
el siguiente bloque cada vez que esto sea posible, y continúa en el
bloque con etiqueta ETI 1 cuando no sea posible. Si no es posible
continuar en ninguno de los bloques, entonces, la transacción espera
hasta que sea posible continuar en alguno de los bloques, dándole
preferencia al siguiente bloque. El último caso indica una transferencia
incondicional al bloque con etiqueta ETI 1.

Utilizando estos últimos bloques, la codificación de este problema se
muestra en la figura 7.5. Note que en esta última codificación se ha
introducido otro estatuto de control llamado STORAGE. En el campo de
localización de este estatuto se identifica al STORAGE y en el campo A
del estatuto se indica la capacidad del STORAGE. De acuerdo a esta
codificación, se obtienen los resultados que se muestran en la figura
7.6. En esta figura se puede observar que cuando hayan pasado
estudiantes por el C.E.C., 290 estudiantes que encontraron ocupadas las
terminales al llegar, ya no regresan. También, se puede apreciar que
280 estudiantes que encontraron al llegar a todas las terminales
ocupadas, sí regresaron.

Por otra parte, la información que se obtiene con respecto al salón de
terminales es la siguiente: El promedio de terminales ocupadas es
9.252, la cantidad de estudiantes que entraron al salón fue de 1009, el
tiempo promedio de utilización de las terminales fue de 59.311 minutos,
el porcentaje de utilización de las terminales es de 92.5\%,
actualmente existen 9 estudiantes trabajando en las terminales y la
cantidad máxima de estudiantes que hubo en el salón fue de 10.

\begin{table}[H]
\centering
\caption*{\textbf{Figura 7.5} Codificación del problema del centro electrónico de cálculo}
\begin{tabular}{lll}
\toprule
 & Generate & 5, 3 \quad \textit{Llegadas al C.E.C.} \\
OTR & Transfer & Both,, LLE \quad \textit{Ver si hay terminales disponibles} \\
 & Enter & 1 \quad \textit{Entrar al salón} \\
 & Advance & 60, 20 \quad \textit{Ocupar la terminal} \\
 & Leave & 1 \quad \textit{Salir del salón} \\
 & Tabulate & 1 \quad \textit{Tiempo en las terminales} \\
 & Terminate & 1 \quad \textit{Se incrementa contador} \\
LLE & Transfer & .5,, REJ \quad \textit{50\% regresa de nuevo} \\
 & Seize & 1 \quad \textit{Entra a la cafetería} \\
 & Advance & 5 \quad \textit{Compra golosinas} \\
 & Release & 1 \quad \textit{Sale de la cafetería} \\
 & Transfer & ,OTR \quad \textit{Regresa al C.E.C.} \\
REJ & Terminate & \quad \textit{Abandonan el C.E.C.} \\
1 & Storage & 10 \quad \textit{Salón con 10 terminales} \\
1 & Table & M1,40,5,50 \quad \textit{Histograma del tiempo en el C.E.C.} \\
 & Start & 10, NP \quad \textit{Elimina estado transiente} \\
 & Reset & \\
 & Start & 1000 \quad \textit{Salida de 1000 clientes} \\
 & End & \\
\bottomrule
\end{tabular}
\end{table}

Por último, la figura 7.6 muestra el histograma de frecuencia del
tiempo de permanencia en el salón de terminales. Como se puede observar
en esta figura, el tiempo promedio que un estudiante se pasa en el
salón de terminales es de 61.047 minutos.

\subsection*{Ejemplo 7.3. Peluquería con tiempo entre llegadas y tiempo de servicio exponenciales}

Para el ejemplo 7.1, suponga que el tiempo entre llegadas de los dos
tipos de clientes es exponencial con media de 65 minutos para el primer
tipo de cliente y con media de 75 minutos para el segundo tipo de
cliente. También, en este mismo ejemplo piense que el tiempo de servicio
es exponencial con media de 15 minutos para un corte de pelo y 7 minutos
para una rasurada. Por último, imagine que en la peluquería solamente
existen 3 lugares disponibles para esperar el servicio.

Para resolver este problema, es necesario introducir una serie de
conceptos nuevos. Por ejemplo, se requiere conocer la forma de
representar dentro del programa una función determinada, así como
también la forma de simular el tiempo entre llegadas y el tiempo de
servicio de una función exponencial.

En GPSS es posible incluir dentro del programa una serie de funciones.
Cada función se define por medio de parejas coordenadas $(X,Y)$. Las
funciones que se manejan pueden ser continuas o discretas. Sin embargo,
puesto que en GPSS las funciones se definen por medio de parejas
coordenadas, entonces, cada vez que se trabaja con funciones continuas,
se interpola linealmente entre las parejas coordenadas correspondientes.
Lo anterior significa que habrá menos error entre más cercanos estén los
puntos que definen la función continua. Por otra parte, no existe
ninguna restricción con respecto a los valores de $X$ y $Y$, a
excepción de $X$ que en algunas ocasiones representa a un número
uniforme entre 0 y 1. Finalmente, conviene señalar que si se referencia
a un valor de $X$ menor que el más pequeño de las parejas coordenadas
que definen la función $((X_1,Y_1),(X_2,Y_2),\ldots,(X_n,Y_n))$,
entonces, el valor de $Y$ que resulta es $Y_1$ y si se referencia un
valor de $X$ mayor que $X_n$, entonces se obtiene para $Y$ un valor de
$Y_n$.

Para hacer referencia a una función, se utiliza el siguiente estatuto de
control:

\begin{center}
ETI FUNCTION A, B
\end{center}

donde:
\begin{align*}
\text{ETI} &= \text{Etiqueta en el campo de localización que identifica a la función.} \\
\text{Campo A} &= \text{Variable independiente que puede ser un gobernador de números uniformes o cualquier otro parámetro} \\
\text{Campo B} &= \text{Identifica si la función que se está usando es continua o discreta. Por ejemplo, } C_n \text{ representa a una función continua que está definida por } n \text{ parejas coordenadas y } D_n \text{ representa a una función discreta que está definida por } n \text{ parejas coordenadas.}
\end{align*}

Inmediatamente después del estatuto de control que hace referencia a la
función, se alimentan las parejas coordenadas que se seleccionaron para
definir la función. Estas parejas se alimentan a partir de la columna 1,
separando con comas los valores de $X$ de los valores de $Y$, y
separando con una diagonal las parejas coordenadas. También, conviene
aclarar que los valores alimentados de $(X,Y)$ no se deben pasar de la
columna 71.

Puesto que el ejemplo que se está analizando, la distribución del tiempo
entre llegadas y el tiempo de servicio es exponencial, entonces, de
acuerdo a la teoría presentada en el capítulo 4 (Método de la
transformada inversa), las parejas coordenadas que se deben alimentar
corresponden a la siguiente función:

\begin{equation}
Y = -\ln(1-R)
\end{equation}

es decir, el valor de $X$ se obtiene por medio de un generador de
números rectangulares\footnote{Un generador de números rectangulares se
presenta como $RN_n$ donde $n$ indica el número del generador que se
está utilizando.} y el valor de $Y$ de acuerdo a la expresión anterior.
Conociendo la media de la distribución exponencial y el valor de $Y$,
entonces, es posible obtener un valor simulado de esa distribución, esto
es:

\begin{equation}
\text{Valor simulado de dist. exponencial} = (\text{Medida de dist. exponencial})(Y)
\end{equation}

este procedimiento es muy recomendable puesto que de esta forma
distribuciones exponenciales con diferentes medias pueden ser simuladas
dentro del mismo programa.

Por otra parte, para simular el tiempo entre llegadas y el tiempo de
servicio, es necesario introducir una variante de los bloques GENERATE
y ADVANCE. Para utilizar estos bloques con distribuciones exponenciales,
el campo A de estos bloques representa la media de la distribución y el
campo B hace referencia a la función $Ln(1-R)$. Conociendo la media y un
valor simulado de $Y$, es posible obtener un valor simulado para el
tiempo entre llegadas o el tiempo de servicio.

Utilizando estos últimos conceptos, la codificación de este problema es
como se muestra en la figura 7.7. De acuerdo a esta codificación, se
obtienen los resultados que se muestran en la figura 7.8. En esta figura
se puede apreciar que por el sistema pasaron 537 clientes del primer
tipo y 463 clientes del segundo tipo. También en esta figura se puede
observar que hubo 63 clientes que no encontraron espacio disponible al
momento de arribar a la peluquería. Por otra parte, el tiempo promedio
de servicio (ver FACILITIES) sin tomar en cuenta al tipo de cliente es
de 17.708 minutos. Además, en esta ocasión el servidor se encuentra
ocupado en 49.5\% de su tiempo disponible.

Con respecto a la sala de espera, la información obtenida es la
siguiente: el contenido promedio es de 0.763, al salón de espera
entraron 1000 clientes, el tiempo promedio de ocupación de un lugar de
la sala de espera fue de 27.305 minutos, la ocupación promedio de un
lugar de la sala de espera fue de 25.4\% y la cantidad máxima de
clientes que hubo en el salón de espera fue de 3.

Finalmente, en la figura 7.8 se muestran los histogramas de los tiempos
de permanencia en la peluquería de los dos tipos de clientes. Como se
puede observar en esta figura, el primer tipo de cliente permanece en
el sistema un tiempo promedio de 24.247 minutos y el segundo tipo de
cliente permanece 30.85 minutos.

\begin{table}[H]
\centering
\caption*{\textbf{Figura 7.7} Codificación del ejemplo 7.3 al considerar tiempo entre llegadas y tiempo de servicio exponencial}
\begin{tabular}{lll}
\toprule
\multicolumn{3}{l}{1 \quad Function \quad RN1, C24} \\
\midrule
\multicolumn{3}{l}{0.0,0.0/0.1,0.104/0.2,0.222/0.3,0.355/0.4,0.509/0.5,0.69} \\
\multicolumn{3}{l}{0.6,0.915/0.7,1.2/0.75,1.38/0.8,1.6/0.84,1.83/0.88,2.12} \\
\multicolumn{3}{l}{0.9,2.3/0.92,2.52/0.94,2.81/0.95,2.99/0.96,3.2/0.97,3.5} \\
\multicolumn{3}{l}{0.98,3.9/0.99,4.6/0.995,5.3/0.998,6.2/0.999,7/0.9997,8} \\
\midrule
 & Generate & 65, FN1 \quad \textit{Llegadas del 1er. tipo de cliente} \\
 & Transfer & Both,,Away \quad \textit{Ver si hay lugar disponible} \\
 & Enter & 1 \quad \textit{Entra a la sala de espera} \\
 & Seize & 1 \quad \textit{Se inicia el servicio} \\
 & Advance & 15, FN1 \quad \textit{Corte de pelo} \\
 & Release & 1 \quad \textit{Se termina el servicio} \\
 & Tabulate & 1 \quad \textit{Tiempo en el sis. del 1er. tipo de cliente} \\
 & Leave & 1 \quad \textit{Deja un espacio en la sala de espera} \\
 & Terminate & 1 \quad \textit{Se incrementa contador} \\
 & Generate & 75, FN1 \quad \textit{Llegada del 2do. tipo de cliente} \\
 & Transfer & Both,,Away \quad \textit{Ver si hay lugar disponible} \\
 & Enter & 1 \quad \textit{Entra a la sala de espera} \\
 & Seize & 1 \quad \textit{Se inicia el servicio} \\
 & Advance & 7, FN1 \quad \textit{Servicio de rasura} \\
 & Advance & 15, FN1 \quad \textit{Corte de pelo} \\
 & Release & 1 \quad \textit{Se termina el servicio} \\
 & Tabulate & 2 \quad \textit{Tiempo en el sist. 2do. tipo de cliente} \\
 & Leave & 1 \quad \textit{Deja un espacio en la sala de espera} \\
 & Terminate & 1 \quad \textit{Se incrementa contador} \\
Away & Terminate & \textit{Los que no encuentran lugar} \\
1 & Storage & 3 \quad \textit{Tres lugares para esperar} \\
1 & Table & M1,5,10,20 \quad \textit{Histograma de frecuencias} \\
2 & Table & M1,5,10,20 \quad \textit{Histograma de frecuencias} \\
 & Start & 10, NP \quad \textit{Elimina estado transiente} \\
 & Reset & \\
 & Start & \\
 & End & 1000 \quad \textit{Salida de 1000 clientes} \\
\bottomrule
\end{tabular}
\end{table}

%==========================================================
\section*{PROBLEMAS}

\begin{enumerate}
\item[7.1.] Cierta máquina produce partes a una razón de 1 pieza cada
15 minutos. El tiempo requerido para inspeccionar estas piezas sigue
una distribución uniforme entre 8 y 14 minutos. El inspector que
examina las piezas aproximadamente acepta el 90\% de las piezas
examinadas. Sin embargo, si una pieza es rechazada, ésta pasa a ser
reprocesada. El tiempo de reproceso sigue una distribución uniforme
entre 10 y 18 minutos. Si se simula la inspección de 1000 piezas,
¿Cuál es el tiempo promedio que una pieza permanece en el sistema?

\item[7.2.] Los estudiantes del ITESM llegan a la copiadora xerox de la
biblioteca a una razón de uno cada $10 \pm 15$ minutos (distribución
uniforme). El tiempo de servicio sigue una distribución uniforme entre
6 y 10 minutos. Sin embargo, si el estudiante encuentra en la cola más
de 10 personas esperando, él decide dar un recorrido de 10 minutos por
la biblioteca. Si se simula el servicio de 1000 estudiantes, determine
el contenido promedio de la cola y el tiempo promedio de permanencia en
el sistema.

\item[7.3.] Al comedor de estudiantes del ITESM llegan dos tipos de
clientes. El tiempo entre llegadas de los dos tipos de clientes es
uniforme entre 5 y 11 minutos. El 65\% de las personas que llegan al
comedor son profesores y el resto estudiantes. El tiempo requerido para
servir los alimentos a un profesor es uniforme entre 2 y 6 minutos, y el
tiempo de servicio para un estudiante es uniforme entre 3 y 7 minutos.
Si la capacidad del comedor es de 200 personas, determine el tiempo
promedio que cada tipo de cliente permanece en el sistema, al momento
de que hayan pasado por el comedor 2000 personas.

\item[7.4.] Resolver utilizando GPSS el problema 5.12.
\item[7.5.] Resolver utilizando GPSS el problema 5.13.
\item[7.6.] Para el problema 5.14, determine mediante GPSS la cantidad
de clientes que al llegar no encuentran lugar en el estacionamiento.
Resuelva este mismo problema suponiendo que el estacionamiento tiene 4,
5, 7 y 8 lugares disponibles.
\item[7.7.] Resolver utilizando GPSS el problema 5.16.
\item[7.8.] Resolver utilizando GPSS el problema 5.17.
\item[7.9.] Los estudiantes del ITESM arriban a la biblioteca de
acuerdo a una distribución exponencial con media de 3 minutos. La
biblioteca consta de 4 pisos con capacidades de 250 para \ldots
\emph{(problema incompleto en el original consultado)}
\end{enumerate}

\end{document}
